\hmsubsection{Team Organization}{OAH}{}

The project group consists of five members with complementary academic backgrounds, 
as presented in Table~1. Four members specialise in Computer Engineering 
(Cyber-Physical Systems), while one member specialises in Mechanical Engineering 
(Product Development). This interdisciplinary composition was essential for developing 
a complete system that integrates mechanical design, embedded systems, and computer 
vision software.

The team is organized with clearly defined roles and associated responsibilities. 
The Team Leader was responsible for overall coordination and project progress. 
The System Lead was responsible for system architecture and integration between 
software and hardware components. The Test Lead was responsible for planning and 
executing testing and system verification. The Documentation \& Scrum Master was 
responsible for maintaining project documentation and facilitating Scrum processes, 
including sprint planning, stand-ups, and retrospectives. The Mechanical Design 
responsible led the physical design of the robot arm, including structural design, 
material selection, and physical implementation.

Although roles were clearly defined, the team worked iteratively and collaboratively 
throughout the project, with tasks distributed based on both expertise and current 
project needs. Communication was maintained through regular internal meetings, weekly 
client meetings every Monday, and supervisor meetings every Thursday. These interactions 
ensured continuous progress monitoring, risk management, and alignment with 
project objectives.